% 08-appendix.tex — 附录

\appendix

\section{完整公式}
\label{sec:appendix-a}

\subsection{FamilyMask 染色体判定 (完整形式)}

设第 $i$ 族 $\mathcal{F}_i = (\text{cells}_i, \text{birth}_i, \text{survive}_i, \text{priority}_i)$, 全局 active 集合 $A \subseteq \{0, \ldots, 15\}$。对元胞 $(x, y)$ 在时刻 $t$ 的更新:
\begin{align*}
  n_i(x, y) &= \sum_{(\Delta x, \Delta y) \in \text{cells}_i} \mathbb{1}\!\left[G_t(x + \Delta x, y + \Delta y) > 0\right] \\
  i^* &= \min\bigl\{\,i \in A : (G_t(x,y) = 0 \wedge n_i(x,y) \in \text{birth}_i) \vee (G_t(x,y) > 0 \wedge n_i(x,y) \in \text{survive}_i)\,\bigr\} \\
  G_{t+1}(x, y) &=
  \begin{cases}
    1, & G_t(x, y) = 0 \wedge n_{i^*}(x, y) \in \text{birth}_{i^*} \\
    1, & G_t(x, y) > 0 \wedge n_{i^*}(x, y) \in \text{survive}_{i^*} \\
    0, & \text{otherwise}
  \end{cases}
\end{align*}

若 $i^*$ 不存在 (即没有 active 族匹配), 元胞 $(x, y)$ 退化为 $\texttt{defaultState} = 0$。

\subsection{maze\_quality 子度量推导 (展开形式)}

记 $V = \{v : G(v) = 1\}$, $|V| = N$, $\deg(v) = \#\{u \sim v : G(u) > 0\}$, $L_{\max}$ 为 4 邻域下的最长简单路径, $C_{\max}$ 为最大 4 连通分量, $a_{\text{outer}}$ 为四条边上的活格子总数 (4 角去重), $H(\cdot)$ 为 Shannon 熵 (自然对数除以 $\ln 2$ 给出比特单位)。
\begin{align*}
  M_b &= \frac{\mathrm{clip}\!\left(\#\{v \in V : \deg(v) = 2\} / N - 0.4,\, 0,\, 0.5\right)}{0.5} \\
  M_s &= \begin{cases}
    \min\!\left(1,\, \dfrac{L_{\max}/N}{0.3}\right), & \dfrac{L_{\max}}{N} \le 0.5 \\
    \max\!\left(0,\, 1 - \dfrac{L_{\max}/N - 0.5}{0.5}\right), & \dfrac{L_{\max}}{N} > 0.5
  \end{cases} \\
  M_j &= \exp\!\left[-\left(\frac{\#\{v : \deg(v) \ge 3\} / N - 0.10}{0.10}\right)^{2}\right] \\
  M_c &= \mathrm{clip}\!\left(\frac{|C_{\max}|}{0.8 N},\, 0,\, 1\right) \\
  M_{\text{bnd}} &= \mathrm{clip}\!\left(1 - \frac{a_{\text{outer}} - 4}{76},\, 0,\, 1\right) \\
  M_p &= \frac{H(\text{2x2 patch})}{\log_2 16} \\
  M_a &= \begin{cases}
    0.5, & \max \mathrm{match}(G, G^{\text{flip}}) < 0.55 \\
    1.0, & 0.55 \le \max \mathrm{match} < 0.85 \\
    1 - \dfrac{\max \mathrm{match} - 0.85}{0.15}, & \max \mathrm{match} \ge 0.85
  \end{cases} \\
  M_t &= \frac{H(\text{2x2 pair})}{\log_2 256} \\
  M_{\text{wall\_gate}} &=
  \begin{cases}
    0, & M_{\text{wall}} \in \{0, 1\} \\
    M_{\text{wall}} / 0.4, & M_{\text{wall}} \in [0, 0.4] \\
    1.0, & M_{\text{wall}} \in [0.4, 0.6] \\
    (1 - M_{\text{wall}}) / 0.4, & M_{\text{wall}} \in [0.6, 1]
  \end{cases}
\end{align*}

其中 $M_{\text{wall}} = \#\{v : G(v) = 0\} / (W H)$ 为墙比。代码实现见 \texttt{src/metrics/maze\_quality.js}。

\subsection{Bellot 2021 $F = \nu/\delta$}

Bellot 2021 提出 $F(M) = \nu(M) / \delta(M)$ (Bellot 2021 §3.4), 本文对照基线使用 \texttt{src/gpu/bellot\_metrics.js::bellotF} 真实实现, 其中 $\nu$ = 非显著墙计数 (死端邻接类), $\delta$ = twistiness proxy (diameter / diagonal)。数值方向: 真迷宫 $F$ 小, 伪迷宫 $F$ 大 (与 $maze\_quality$ 相反, 详见 §4.5)。

\section{数据可用性}
\label{sec:appendix-b}

所有产物在 GitHub 与本地开源:
\begin{itemize}[leftmargin=*, itemsep=2pt]
  \item \textbf{项目源码}: \texttt{E:/doro/maze-web/} \\
    \texttt{node index.html} 启动 (默认 \texttt{python -m http.server 8087} 提供静态托管)。
  \item \textbf{小范围 sweep ndjson}: \texttt{sweep\_2026\_07\_04/results.ndjson} \\
    含 122 OK + 6 TIMEOUT, 每行含 mask / maxFam / seed / best / wall\_sec / status。
  \item \textbf{大范围 sweep ndjson}: \texttt{sweep\_2026\_07\_08\_big/results.ndjson} \\
    含 5/5 OK, manhattan-2 / maxFam=8 × \{s444, s1111, s2222, s3333, s5555\}。
  \item \textbf{15-pattern 图案数据}: \texttt{paper/data/\_grids.json} \\
    15 张 $31 \times 31$ 网格, 0/1 flat 数组, 总 15 $\times$ 961 = 14415 个元素。
  \item \textbf{15-pattern 分数}: \texttt{paper/data/sweep\_summary.json::15pattern.per\_pattern} \\
    每个 pattern 含 name, type, mq, wall\_ratio, nu\_count, twistiness, bellotF。
  \item \textbf{15-pattern 复核脚本}: \texttt{paper/data/\_eval\_15patterns.mjs} \\
    用 \texttt{src/metrics/maze\_quality.js::mazeQuality} 与 \texttt{src/gpu/bellot\_metrics.js::bellotF} 真实源码重算 15 个 pattern, 与 \texttt{per\_pattern} 中缓存值对比。复核结果: 15/15 mq 精确匹配, 15/15 Bellot $F$ 精确匹配 (容差 mq $\le 0.005$, $F$ $\le 0.05$)。复核输出保存于 \texttt{paper/data/\_15pat\_recompute.json}。
  \item \textbf{示意图生成脚本}: \texttt{tools/regen\_fig\_15pat\_grids.py}, \texttt{tools/regen\_fig\_sweep\_heatmaps.py}。
\end{itemize}

\section{硬件与软件环境}
\label{sec:appendix-c}

\begin{itemize}[leftmargin=*, itemsep=2pt]
  \item \textbf{Host}: Windows 11 (x64), \texttt{python=3.13.0}, \texttt{uv} 包管理。
  \item \textbf{Browser}: Chromium 114+ (WebGPU compute shader), 实测 Edge 119.0.2151.72。
  \item \textbf{CA/GPU kernel}: WGSL via \texttt{src/gpu/*.js}, dispatch 维度 $(\text{ruleCount} \times \text{seedCount}, \text{cellCount})$。
  \item \textbf{小范围 sweep}: pop=200, gens=500, $40 \times 60$ 网格, 300 步 CA; $\sim$ 12.5 小时 128 次 (122 OK + 6 TIMEOUT), 单 run $\sim$ 4 分钟。
  \item \textbf{大范围 sweep}: pop=500, gens=2000, $500 \times 2000$ 网格, 300 步 CA; $\sim$ 4.2 小时 5/5 run, 单 run $\sim$ 50 分钟。
  \item \textbf{Paper tooling}: \texttt{xelatex} (TeX Live 2024), \texttt{matplotlib} 3.8 (单文件 \texttt{.png} 输出), Python 脚本位于 \texttt{tools/}。
\end{itemize}

\begin{thebibliography}{99}
\small
\bibitem{conway1970} Conway, J. (1970). \textit{The Game of Life}. Scientific American, 223(4), 4--7.
\bibitem{wolfram2002} Wolfram, S. (2002). \textit{A New Kind of Science}. Wolfram Media, Champaign IL.
\bibitem{mitchell1993} Mitchell, M., Hraber, P., \& Crutchfield, J. (1993). Revisiting the Edge of Chaos. \textit{Complex Systems}, 7, 89--130.
\bibitem{mordvintsev2020} Mordvintsev, A., Randazzo, E., Niklasson, E., \& Levin, M. (2020). Growing Neural Cellular Automata. \textit{Distill}, 5(8), e24.
\bibitem{palm2022} Palm, R., et al. (2022). Variational Neural Cellular Automata. \textit{ICLR 2022}.
\bibitem{bellot2021} Bellot, V., Cautrès, M., Favreau, J.-M., Gonzalez-Thauvin, M., Lafourcade, P., Le Cornec, K., Mosnier, B., \& Rivière-Wekstein, S. (2021). How to generate perfect mazes? \textit{Information Sciences}, 572, 444--459. \href{https://doi.org/10.1016/j.ins.2021.03.022}{doi:10.1016/j.ins.2021.03.022}
\bibitem{pech2002} Pech, J. (2002). Perfect Maze Generation. Master's thesis, Charles University Prague.
\bibitem{rejbrand2017} Rejbrand, A. (2017). Cellular automaton maze generation. \textit{Andreas Rejbrand's Website}, March 5, 2017. \href{https://english.rejbrand.se/rejbrand/article.asp?ItemIndex=421}{URL}
\bibitem{mcclendon2001} McClendon, M. S. (2001). The complexity and difficulty of a maze. In \textit{Proceedings of Bridges 2001: Mathematical Connections in Art, Music, and Science}, Winfield, Kansas, July 27--29, 2001, pp. 213--222.
\bibitem{liu2018darts} Liu, H., Simonyan, K., \& Yang, Y. (2018). DARTS: Differentiable Architecture Search. \textit{ICLR 2019}.
\bibitem{real2019regularized} Real, E., Aggarwal, A., Huang, Y., \& Le, Q. (2019). Regularized Evolution for Image Classifier Architecture Search. \textit{AAAI 2019}.
\bibitem{webgpu2023} W3C / GPU for the Web Working Group (2023). WebGPU compute shaders specification, W3C TR.
\bibitem{skiena2024} Skiena, S. (2024). Codeforces with WebGPU. \textit{Codeforces Blog}, 2024.
\bibitem{margolus1984} Margolus, N. (1984). Physics-like model of computation. \textit{Physica D}, 10(1-2), 81--95.
\end{thebibliography}
